\section{Contract Theory}

\subsection{Hidden Action (Moral Hazard)}
Principal offers wage scheme \(w(y)\); agent picks unobserved effort \(e \in \{L, H\}\) at cost \(c(e)\); output \(y \sim f(\cdot|e)\). Agent utility \(u(w) - c(e)\), \(u\) concave; reservation utility \(\bar u\).

\textbf{Program} (implement \(H\)): \(\min_{w(\cdot)} \mathbb{E}[w(y)|H]\) subject to
\begin{gather*}
\text{(IR)} \quad \mathbb{E}[u(w)|H] - c(H) \geq \bar u \\
\text{(IC)} \quad \mathbb{E}[u(w)|H] - c(H) \geq \mathbb{E}[u(w)|L] - c(L)
\end{gather*}

\textbf{First best} (observable \(e\)): IC drops, full insurance---flat wage, effort dictated.

\textbf{Second best:} optimal wage satisfies
\[\frac{1}{u'(w(y))} = \lambda + \mu\left(1 - \frac{f(y|L)}{f(y|H)}\right)\]
Pay is high where the \textbf{likelihood ratio} \(f(y|H)/f(y|L)\) is high---wages reward \textit{evidence of effort}, not output per se. Agent bears risk \(\implies\) implementation costs exceed first best (risk premium is the agency cost).

\textbf{MLRP} (monotone likelihood ratio) \(\implies\) wage increasing in output. \textbf{Informativeness:} condition pay on any signal that adds information about \(e\); filter out common shocks (relative performance).

\subsection{Hidden Information (Screening)}
Principal designs a menu; agent has private type \(\theta \in \{\theta_L, \theta_H\}\) (prob.\ \(p_L, p_H\)). By the \textbf{revelation principle}, restrict to direct menus \(\{(q_L, t_L), (q_H, t_H)\}\) satisfying IC and IR per type.

\textbf{Solution structure} (monopolist screening a buyer):
\begin{itemize}[nosep]
    \item Binding constraints: \textbf{IR of the low type}, \textbf{IC of the high type} (high type is tempted to mimic down); the other two are slack
    \item \textbf{No distortion at the top:} \(q_H\) is efficient
    \item \textbf{Distortion at the bottom:} \(q_L\) below efficient---sacrificing low-type surplus cheapens the high type's \textbf{information rent} \(U_H = \Delta\theta\, q_L\)
    \item Low type may be excluded entirely if \(p_H\) is large
\end{itemize}

\textbf{Same anatomy everywhere:} second-degree price discrimination, optimal income taxation (Mirrlees), regulation of an unknown-cost firm (Baron--Myerson), insurance menus (deductibles screen risk types).

\subsection{Moral Hazard vs.\ Screening: Exam Diagnostics}
\begin{itemize}[nosep]
    \item Private information \textit{before} contracting \(\implies\) screening/selection; \textit{after} \(\implies\) moral hazard
    \item Screening distorts \textbf{quantities} (menus, rationing); moral hazard distorts \textbf{risk-sharing} (pay-performance)
    \item Both leave rents: information rent (screening) vs.\ risk premium (moral hazard)
\end{itemize}
